\section{Define one accepted representation for each mapped value}
Use the shortest byte sequence. For more than one byte, the final payload must be
nonzero. Thus \texttt{01} encodes one, while \texttt{81 00} is rejected even though
its base-128 arithmetic also yields one. Canonicality is a format policy here; do not
assume that every protocol using LEB128 rejects longer equivalent encodings.

A 33-bit unsigned value needs at most five bytes. Four groups account for 28 bits,
leaving five usable payload bits in the last possible byte. At byte five, require
continuation zero and payload less than 32. Payload bits five and six must be zero.
Check this before reconstructing a reading. A fifth payload of 32 is a malformed
wire value, even if a later range check would also reject the resulting number.

\begin{table}[htbp]
\centering
\caption{Wire validity is distinct from the reconstructed reading range.}
\tagpdfsetup{table/header-rows={1}}
\begin{tabular}{ll}
\toprule
Bytes & Interpretation at previous reading zero\\
\midrule
\texttt{00} & Valid delta zero\\
\texttt{80} & Malformed: missing terminating byte\\
\texttt{80 00} & Malformed: nonminimal zero\\
\texttt{80 80 80 80 20} & Malformed: fifth payload exceeds five bits\\
\texttt{80 80 80 80 80 01} & Malformed: byte five still continues\\
\texttt{ff ff ff ff 1f} & Canonical 33-bit value; reconstructed reading out of range\\
\bottomrule
\end{tabular}
\end{table}

The first five-byte encoder examples expose the width requirement:
\texttt{ff ff ff ff 0f} represents an initial minimum reading, and
\texttt{fe ff ff ff 0f} represents an initial maximum reading.
The upward endpoint jump uses \texttt{fe ff ff ff 1f}; the downward jump uses
\texttt{fd ff ff ff 1f}. A parser that permits only four useful bits in byte five
would incorrectly reject legitimate endpoint transitions.

\begin{MARSExercise}{Locate each rejection}
For every malformed row in the table, identify the earliest byte at which the parser
can decide it is invalid. For the truncated single-byte row, identify the input-length
comparison that supplies the missing evidence. Do not read a sixth byte to reject
continuation in the fifth.
\end{MARSExercise}
